\subsection{Gauge-chain status and falsification boundary}

The finite carrier calculation derives the incidence polynomial for \(J\),
the target-blind response \(R=-J\), its relative sector signs, and the
equivariant current algebra
\(\mathfrak{su}(3)\oplus\mathfrak{su}(2)\oplus\mathfrak u(1)\).
Given the declared fermionic Spin category and rank-15 matter projector pair,
the matter certificate derives anomaly balance, exact hypercharge up to
conjugation, three colors, the common \(\mathbb Z_6\) tensor kernel, and the
maximal faithful matter image. These implications do not use Minimal
Admissible Realization.

Four central global forms act consistently on the same local tensors: the
cover and its \(\mathbb Z_2\), \(\mathbb Z_3\), and \(\mathbb Z_6\)
quotients. The kernel theorem identifies the maximal faithful local matter
image; it does not select the physical spectrum of line operators or bundles.
The physical global form is open. Physical fermion typing,
laboratory-current attachment, equality between the finite current and the
categorically reconstructed group, scalar attachment and multiplicity,
physical family attachment, and QFT realization are also open.

The gauge result fails if the incidence identity, target-blind response,
equivariant lift, anomaly scan, tensor-character enumeration, or refinement
persistence receipt fails its registered negative controls. The categorical
route fails outside its combined strictification, trivial-holonomy,
localization, and refinement/fiber receipts.

\begin{center}\rule{0.5\linewidth}{0.5pt}\end{center}
